\begin{definition}[Finite counterterms along $\Delta$ and $P_b$]
\label{def:c11-free}
Write the solutions of the arity-one equations along $\Delta$ and along $P_b$ as
\[
\begin{aligned}
\hbar F_1^{(1)}(\Delta)=\ &4\,\text{units}\cdot(\text{field-counting functional})+\Big[-\tfrac23\int(A,A^*)+\tfrac13\int(\lambda^*,\lambda)+\tfrac43\int(c^*,c)\Big]\,\text{units}\\
&+\theta\,\text{units}\int\big(2E(A|_{W_{\mathbb R}^\perp})+2A|_{W_{\mathbb R}^\perp},A^*|_{W_{\mathbb R}^\perp}\big)-\theta_\lambda\,\text{units}\int\big(2E(\lambda)+3\lambda,\lambda^*\big)\\
&+\theta_V\,\text{units}\Big[\int\big(2E(A|_V)+2A|_V,A^*|_V\big)-\int\big(2E(c),c^*\big)\Big]+\text{units}\cdot Q\Big(\kappa_V\int(\iota_xA|_V,c^*)+\kappa_\lambda\int(\rho(x)\lambda^*,\lambda^*)\Big)+(\text{cubic and higher}),\\
\hbar F_1^{(1)}(P_b)=\ &\theta_P\,\text{units}\int\big(\partial_b(A|_{W_{\mathbb R}^\perp}),A^*|_{W_{\mathbb R}^\perp}\big)-\theta_{P,\lambda}\,\text{units}\int(\partial_b\lambda,\lambda^*)+\theta_{P,V}\,\text{units}\Big[\int\big(\partial_b(A|_V),A^*|_V\big)-\int(\partial_bc,c^*)\Big]\\
&+Q\big(\text{a combination of }\textstyle\int(\iota_bA|_V,c^*)\text{ and }\int(\rho(b)\lambda^*,\lambda^*)\big)+(\text{cubic and higher}),
\end{aligned}
\]
with numbers $\theta,\theta_\lambda,\theta_V,\kappa_V,\kappa_\lambda,\theta_P,\theta_{P,\lambda},\theta_{P,V}$; $\theta$ is the coefficient of \eqref{eq:c9-Fperp} and $\theta_P$ that of Theorem \ref{thm:c9-P1}. The three terms with $\theta$, $\theta_\lambda$, $\theta_V$ are the sector parts of $S^{(1)}(\Delta)$ in the sectors of $A|_{W_{\mathbb R}^\perp}$, of $\lambda$ and of $(A|_V,c)$ times these numbers, and those with $\theta_P$, $\theta_{P,\lambda}$, $\theta_{P,V}$ the sector parts of $S^{(1)}(P_b)$ (Definition \ref{def:c11-kernels}). The scheme is translation covariant if the position part along $b$ of $\hbar F_1^{(1)}(\Delta)$ is $2\hbar F_1^{(1)}(P_b)$; then $\theta_P=\theta$, $\theta_{P,\lambda}=\theta_\lambda$, $\theta_{P,V}=\theta_V$. The arity-two equations along $(\Delta,P_b)$ and $(\Delta,\xi)$ fix $\theta=\theta_\lambda=\theta_V=-\tfrac13$ (Theorem \ref{thm:c9-theta2}), and those along $(\xi,P_b)$ give $\theta_P=\theta_{P,\lambda}=\theta_{P,V}$ (Theorem \ref{thm:c11-oddar2}).
\end{definition}
